\documentclass[11pt]{article}
\usepackage[margin=1in]{geometry}
\usepackage{amsmath,amssymb,amsthm}
\usepackage{enumitem}

\newtheorem{theorem}{Theorem}[section]
\newtheorem{lemma}[theorem]{Lemma}
\newtheorem{remark}[theorem]{Remark}

\newcommand{\dT}{\delta_T}
\newcommand{\al}{\alpha}
\newcommand{\Et}{E_{\hat t}}
\newcommand{\dH}{\,d\mathcal H^{N-1}}

\title{Assessment of a Selection--Plan 3 Hybrid Graph Entry}
\author{}
\date{May 2026}

\begin{document}
\maketitle

\section{The Proposed Replacement}

The proposed route is:
\[
\hbox{selection}
\Rightarrow
\hbox{controlled class}
\Rightarrow
\hbox{Plan 3 good inset level } \Et
\Rightarrow
L^2\hbox{-Serrin}
\Rightarrow
L^\infty\hbox{-Serrin}
\Rightarrow
\hbox{Brandolini annulus}
\Rightarrow
L^\infty\hbox{ radial graph.}
\]
The conclusion of this check is mixed.  The first half is genuinely useful:
selection turns the weighted \(D_H\) variance into honest unweighted \(L^2\)
Serrin control on the stopped level, and the inward cut preserves the bad
quotient/asymmetry.  The weak point is the \(L^2\to L^\infty\) upgrade.  It is
not a consequence of Lipschitz boundary geometry alone.

\section{What Selection Gives}

Let \(U\) be the selected minimizer from Plan 1, with torsion function \(u\).
The BDV/selection regularity package gives a uniform Lipschitz bound for the
potential:
\[
   \|\nabla u\|_{L^\infty(U)}\le L_*(N,R).
   \tag{2.1}\label{eq:lipu}
\]
It also gives density and nondegeneracy.  Before the flatness/graph-entry
step, it does \emph{not} give a global Lipschitz parametrization of
\(\partial U\).  Thus the statement that selection reduces to Lipschitz
boundary is stronger than the current Plan 1 selection lemma.  If one assumes
a uniform Lipschitz boundary class separately, the discussion below is
conditional on that assumption.

\section{Chebyshev Level Gives Honest \(L^2\)}

Apply the Plan 3 volume-variable Chebyshev selection inside \(U\).  For every
fixed small \(\theta>0\), there is a regular level
\[
   E=\Et:=\{u>\hat t\}
\]
such that
\[
   |U\setminus E|\le C_\theta \dT(U),
   \qquad
   D_H(\hat t)\le C\theta.
   \tag{3.1}\label{eq:good-level}
\]
On this level the exact identity gives, with \(f=|\nabla u|\) and
\(\bar f=|E|/P(E)\),
\[
   \int_{\partial E}\frac{(f-\bar f)^2}{f}\,\dH
      =
      \frac{|E|}{P(E)^2}D_H(\hat t).
   \tag{3.2}
\]
Using \eqref{eq:lipu},
\[
\begin{aligned}
   \int_{\partial E}(f-\bar f)^2\,\dH
      &=
      \int_{\partial E}
      f\,\frac{(f-\bar f)^2}{f}\,\dH  \\
      &\le L_*(N,R)
      \int_{\partial E}\frac{(f-\bar f)^2}{f}\,\dH
      \le C(N,R)\theta.
\end{aligned}
\tag{3.3}\label{eq:l2serrin}
\]
This is the positive part of the hybrid route: selection removes the
upper-gradient tail obstruction that raw \(D_H\) had.

\section{The Inward Cut Preserves Asymmetry}

Let
\[
   q:=\frac{E(U)-E(B)}{\alpha(U)}.
\]
Since \(\dT(U)=q\,\alpha(U)\), \eqref{eq:good-level} gives
\[
   |U\setminus E|\le C_\theta q\,\alpha(U).
\]
The smooth asymmetry \(\alpha\) is \(L^1\)-Lipschitz on bounded sets.  After
volume normalization of \(E\),
\[
   \alpha(E^\sharp)
      \ge
      \alpha(U)-C|U\setminus E|-C|\ |E|-|U|\ |
      \ge
      (1-C_\theta q)\alpha(U).
   \tag{4.1}\label{eq:asym-preserve}
\]
Thus in the selection contradiction regime \(q\ll\theta\), the inset level
keeps the asymmetry comparable to that of \(U\).  This part of the proposed
route is sound.

\section{The Missing \(L^2\to L^\infty\) Upgrade}

To use a Brandolini/Serrin annulus theorem one wants
\[
   \|\,|\nabla u|-\bar f\,\|_{L^\infty(\partial E)}
      \le \varepsilon_B
   \tag{5.1}\label{eq:linftyserrin}
\]
with \(\varepsilon_B\ll1\).  The \(L^2\) estimate
\eqref{eq:l2serrin} implies this only if one has a uniform modulus of
continuity for \(f=|\nabla u|\) on \(\partial E\).  For instance, if
\[
   [f]_{C^\gamma(\partial E)}\le M
   \tag{5.2}\label{eq:holder}
\]
and \(\partial E\) has uniformly controlled \((N-1)\)-dimensional geometry,
then interpolation gives
\[
   \|f-\bar f\|_{L^\infty(\partial E)}
      \le
      C\,M^{\frac{N-1}{N-1+2\gamma}}\,
      \|f-\bar f\|_{L^2(\partial E)}^{\frac{2\gamma}{N-1+2\gamma}}.
   \tag{5.3}\label{eq:interp}
\]
Then choosing \(\theta\) small enough would give \eqref{eq:linftyserrin}.

The problem is that \eqref{eq:holder} is not implied by a uniform Lipschitz
boundary.  Lipschitz domains do not give uniform continuity of \(\nabla u\) up
to the boundary.  For the inset level \(\partial E\), qualitative analyticity
is free because the level is inside \(U\), but the quantitative constants
depend on
\[
   \operatorname{dist}(\partial E,\partial U)
   \quad\hbox{and}\quad
   \inf_{\partial E}|\nabla u|.
\]
The stopped level is only \(O_\theta(\dT)\) from the boundary in measure, and
the physical distance can go to zero.  Also \(D_H\) controls the
low-gradient set only in measure:
\[
   \mathcal H^{N-1}\{|\nabla u|\le s\}
      \lesssim D_H\,s,
\]
not pointwise.  Hence it does not give a uniform lower bound for
\(|\nabla u|\) on the entire level.

Therefore the implication
\[
   \hbox{selected Lipschitz class}+D_H(\hat t)\ll1
   \Longrightarrow
   \|\,|\nabla u|-\bar f\,\|_{L^\infty(\partial E)}\ll1
\]
is currently unproved and is false as a general Lipschitz-domain principle.
It becomes true if one adds a uniform \(C^\gamma\) modulus for the boundary
gradient on \(\partial E\), but that is essentially a free-boundary regularity
input.

\section{Inset Boundary Regularity}

For a regular value \(\hat t\), \(\partial E=\{u=\hat t\}\) is a smooth
embedded hypersurface.  This is qualitative.  A quantitative Lipschitz
constant follows from the implicit function theorem with constants depending
on
\[
   \frac{\|D^2u\|_{L^\infty(\hbox{collar})}}
        {\inf_{\partial E}|\nabla u|}.
\]
Neither term is uniformly controlled by \eqref{eq:lipu} and \(D_H(\hat t)\ll1\).
Thus ``the inset boundary is Lipschitz'' is true with constants depending on
the selected level, but not with the uniform constants needed for a proof.

\section{Brandolini Annulus and Graph Entry}

Assume conditionally that \eqref{eq:linftyserrin} is obtained and that
\(\partial E\) lies in the geometric class required by a Brandolini-type
stability theorem.  Then one gets a thin annulus:
\[
   B_{r_i}(z)\subset E\subset B_{r_e}(z),
   \qquad
   r_e-r_i\le \eta_B.
   \tag{7.1}
\]
This gives Hausdorff closeness to a ball.  It does not by itself give a radial
graph.  A set may contain \(B_{r_i}(z)\), be contained in \(B_{r_e}(z)\), and
still have overhangs or extra components in the thin shell.  One still needs a
one-crossing/star-shapedness theorem, or an \(O(\eta_B)\)-cost radialization.

\section{Conditional Version That Would Work}

The proposed route becomes correct under the following additional hypotheses:

\begin{enumerate}[label=\arabic*.]
\item the selected/inset level belongs to a uniform geometric class for which
\(|\nabla u|\) has a \(C^\gamma\) modulus bounded by \(M(N,R)\);
\item Brandolini's annulus theorem applies to that class;
\item thin annulus plus selected topology implies one radial crossing, or an
annulus-to-radialization theorem is available.
\end{enumerate}

Under these assumptions, choose \(\theta\) below the interpolation and
Brandolini thresholds.  Chebyshev gives \eqref{eq:l2serrin}; interpolation
gives \eqref{eq:linftyserrin}; Brandolini gives the annulus; the final
one-crossing/radialization input gives a small \(L^\infty\) radial graph.
Equation \eqref{eq:asym-preserve} and the superlevel torsion transfer preserve
the bad quotient.

\section{Verdict}

The hybrid has a useful core:
\[
   \hbox{selection's Lipschitz bound for }u
   +D_H\hbox{ on a Chebyshev level}
   \Longrightarrow
   L^2\hbox{-Serrin on that level.}
\]
It also preserves asymmetry in the selection contradiction regime.

But it does not yet replace the Alt--Caffarelli graph-entry machinery.  The
serious missing theorem is a uniform \(L^2\to L^\infty\) Serrin upgrade on the
inset level.  Merely assuming Lipschitz boundary is not enough; one needs a
uniform modulus of continuity for the boundary gradient, which is effectively
regularity/flatness information.  A second, smaller missing theorem is the
annulus-to-radial-graph or annulus-to-radialization step.

Thus the best current formulation is not a complete replacement for Plan 1,
but a narrowed target: prove just enough regularity of the \emph{inset level}
to make the interpolation estimate \eqref{eq:interp} valid, instead of proving
global graph regularity of \(\partial U\) directly.

\end{document}
